% =====================================================================
%  Curriculum Vitae -- John "Jack" P. McGuire
%
%  Maintainable LaTeX source recreating the layout of the CV PDF that
%  was previously published to jackmcguireastro.github.io (cv.pdf).
%  No .tex source existed in the repository, so this file was written
%  to match the existing compiled style: 10pt serif, all-caps ruled
%  section headings, bold entry titles with right-aligned dates.
%
%  Build:  pdflatex cv.tex   (run twice so hyperref outlines settle)
%  Output: cv.pdf  -> repository root of the website
%
%  Last content update: August 2026
% =====================================================================
\documentclass[10pt,a4paper]{article}

\usepackage[T1]{fontenc}
\usepackage[utf8]{inputenc}
\usepackage[
  top=0.85in, bottom=0.85in, left=0.85in, right=0.85in
]{geometry}
\usepackage{titlesec}
\usepackage{enumitem}
\usepackage{textcomp}
\usepackage[hidelinks]{hyperref}

\hypersetup{
  pdftitle={Curriculum Vitae -- John P. McGuire},
  pdfauthor={John P. McGuire},
  pdfsubject={Astronomy curriculum vitae},
  pdfkeywords={astronomy, binary evolution, barium dwarfs, interferometry}
}

% ------------------------------------------------------------ spacing
\setlength{\parindent}{0pt}
\setlength{\parskip}{0pt}
\raggedbottom

% ----------------------------------------------------------- headings
% All-caps bold section name with a full-width rule underneath.
\titleformat{\section}
  {\normalfont\normalsize\bfseries}
  {}{0em}{\MakeUppercase}[\vspace{-0.55em}\rule{\textwidth}{0.5pt}]
\titlespacing*{\section}{0pt}{9pt}{3pt}

% ------------------------------------------------------------- macros

% \entry{Left-hand bold title}{Right-aligned date}
\newcommand{\entry}[2]{%
  \par\vspace{2pt}%
  \noindent\textbf{#1}\hfill#2\par
}

% Plain continuation line (e.g. degree, mentor, institution)
\newcommand{\cvline}[1]{\noindent#1\par}

% Compact bullet list
\newenvironment{items}
  {\begin{itemize}[leftmargin=1.2em,itemsep=0.6pt,topsep=1.5pt,parsep=0pt,
                   label=\textbullet]}
  {\end{itemize}}

% "Label: value" line used in the skills section
\newcommand{\skill}[2]{\begingroup\sloppy\noindent\textbf{#1} #2\par\endgroup\vspace{1.5pt}}

% =====================================================================
\begin{document}

% ------------------------------------------------------------- header
\begin{center}
  {\LARGE John ``Jack'' P. McGuire}\\[3pt]
  {\small\href{mailto:jackmcg7@live.com}{jackmcg7@live.com}
   \ $\vert$ \
   \href{https://jackmcguireastro.github.io}{jackmcguireastro.github.io}}
\end{center}
\vspace{4pt}

% ---------------------------------------------------------- education
\section{Education}

\cvline{\textbf{University of Wyoming}, Laramie, WY}
\entry{Ph.D. Student in Physics (Astronomy Concentration)}{August 2026--Present}

\vspace{3pt}
\cvline{\textbf{Georgia State University}, Atlanta, GA}
\begin{items}
  \item \textbf{Bachelor of Science in Physics} \hfill May 2026 \\
        Overall GPA: 3.76; Institutional GPA: 3.81
  \item \textbf{Associate of Science in Physics, High Honors} \hfill May 2024 \\
        Overall GPA: 3.73; Institutional GPA: 3.82
\end{items}

% --------------------------------------------------------- employment
\section{Academic Employment}

\entry{University of Wyoming, Department of Physics \& Astronomy}{August 2026--Present}
\cvline{Graduate Teaching Assistant}
\begin{items}
  \item Support undergraduate physics and astronomy instruction through laboratory or
        recitation teaching, grading, office hours, and individual student assistance.
\end{items}

% ------------------------------------------------------- publications
\section{Publications \& Manuscripts}

\begin{items}
  \item Bentz, M.~C.; Alahakone, R.; Azoulay, L.; Burns-Kaurin, E.; Chaturvedi, A.~S.;
        Davis, M.; Kimani-Stewart, K.; Markham, M.; \textbf{McGuire, J.~P.}; McMullan, E.;
        Patel, M.; Binu Raj, M.; Sankey, D.; Wilson, A.~A.-L.; Zdanky, K. (2025).
        ``Broad-band Monitoring of 797 Montana.'' \textit{Minor Planet Bulletin},
        \textbf{52}(3), 207--208.
  \item \textbf{Where Are the K \& M Barium Dwarfs?} Manuscript in preparation; planned
        submission to an American Astronomical Society journal in Fall 2026.
\end{items}

% ------------------------------------------------------ presentations
\section{Selected Presentations \& Posters}

\begin{items}
  \item \textbf{Finding Barium Dwarfs Among Established White-Dwarf Systems.} Poster
        presented at the Society of Physics Students Research Conference,
        April 23, 2026.
  \item \textbf{Finding Barium Dwarfs Among Established White-Dwarf Systems.} Poster
        presented at the Georgia State Undergraduate Research Conference (GSURC),
        April 6, 2026. Presented spectroscopic abundance-ratio measurements for five
        candidate barium dwarfs.
  \item \textbf{Posters at the Georgia State Capitol.} Selected to represent the Georgia
        State University Department of Physics \& Astronomy and presented research on the
        CHARA Array and the gLOWCOST muon-monitoring network to legislators and the
        public, February 19, 2026.
  \item \textbf{AGB Post-Mass-Transfer Spectroscopic Ratios as White Dwarf Tracers.}
        Poster presented at the Georgia Regional Astronomy Meeting, Emory University,
        November 7--8, 2025.
  \item \textbf{Where Are the K \& M Barium Dwarfs? Potential Answers in the 40 Eridani
        System.} Talk presented at the Georgia State University Undergraduate Research
        Symposium, August 7, 2025.
\end{items}

% ----------------------------------------------------------- research
\section{Research Experience}

\entry{Undergraduate Research Assistant, Georgia State University}{August 2025--August 2026}
\cvline{Mentor: Dr. Russel J. White}
\begin{items}
  \item Investigated the 71 spectroscopically confirmed barium dwarfs using Gaia DR3
        astrometry, radial-velocity statistics, archival spectroscopy,
        spectral-energy-distribution fitting, and MIST evolutionary models.
  \item Developed Python workflows using NumPy, Pandas, Astropy, Astroquery, and
        Matplotlib to cross-match catalogs, distinguish dwarfs from evolved interlopers,
        and compare barium dwarfs with white-dwarf--main-sequence binaries.
  \item Collaborated with Dr. Richard O. Gray of Appalachian State University on a
        high-resolution spectroscopic search for s-process enrichment in GJ~166~C. The
        null result, together with the system's roughly 200-year orbit, constrained the
        separation required for efficient AGB mass transfer.
  \item Continued CHARA interferometric reductions and analysis of nearby binary M dwarfs;
        contributed results to manuscripts on barium dwarfs and M-dwarf radius inflation.
\end{items}

\entry{Raghavan Fellow, Georgia State University Undergraduate Summer Research Program}{May--August 2025}
\cvline{Mentor: Dr. Russel J. White}
\begin{items}
  \item Initiated parallel projects on barium-dwarf binary evolution and CHARA
        interferometric measurements of the GJ~166 system.
  \item Reduced four nights of MIRC-X and MYSTIC data; fitted uniform-disk and
        limb-darkened models, propagated uncertainties, and compared independent
        reduction pipelines.
  \item Identified the strong concentration of known barium dwarfs among F- and G-type
        stars, motivating a search for the missing K- and M-dwarf population.
\end{items}

\entry{Cosmic-Ray Research Assistant, Georgia State University}{2025--2026}
\cvline{Mentors: Dr. Xiaochun He and Dr. Viacheslav Sadykov}
\begin{items}
  \item Assembled Raspberry Pi-based detectors for the gLOWCOST global muon-monitoring
        network and supported detector-based science outreach.
  \item Analyzed multi-year muon-flux measurements using atmospheric data, solar and lunar
        ephemerides, Fourier transforms, coherence analysis, and correlation tests to
        investigate environmental and tidal modulation.
  \item Studied Forbush decreases by comparing gLOWCOST measurements with ACE MAG and
        SWEPAM solar-wind and magnetic-field data from the L1 point.
\end{items}

% ---------------------------------------------------------- observing
\section{Observing Experience}

\entry{Center for High Angular Resolution Astronomy (CHARA) Array}{November 2025}
\begin{items}
  \item Planned and led two half-night observing sessions with PI Becky Flores using the
        six-telescope CHARA Array.
  \item Operated MIRC-X ($H$ band) and MYSTIC ($K$ band) to observe nearby M dwarfs in
        binary systems and measure their angular diameters.
  \item Reduced and modeled archival observations of GJ~166~C, finding an approximately
        30\% radius excess relative to stellar-evolution model predictions.
\end{items}

\entry{Hard Labor Creek Observatory}{February 2025}
\begin{items}
  \item Participated in a four-night photometric campaign on main-belt asteroid
        797 Montana using the 24-inch Miller telescope.
  \item Reduced CCD photometry using IRAF and DS9; the resulting light curve yielded a
        rotation period of $4.545 \pm 0.001$~hr and led to a refereed publication.
\end{items}

% ----------------------------------------------------------- training
\section{Professional Training \& Public Service}

\entry{American Physical Society -- Advocacy Champion}{June 2026--Present}
\begin{items}
  \item Participate in APS science-policy and advocacy training and coordinated outreach
        supporting research funding, STEM education, and the physics workforce.
  \item Build sustained relationships with elected officials and their staff to communicate
        the priorities and public value of the physics community.
\end{items}

\entry{Hello SciCom -- Comedy Camp for Science}{August--October 2026}
\begin{items}
  \item Selected for an eight-week virtual training program in using humor, storytelling,
        and performance to communicate technical STEM research accurately and accessibly.
  \item Develop a three-minute public-facing talk and contribute to a research study
        examining how humor affects trust, engagement, and perceptions of science as a
        public good.
\end{items}

% ------------------------------------------------------------- skills
\section{Skills \& Interests}

\skill{Scientific Computing:}{Python (NumPy, Pandas, Matplotlib, Astropy, Astroquery,
Specutils, Lightkurve), TOPCAT, \LaTeX{}, HTML}

\skill{Astronomy Tools:}{Gaia Archive/ADQL, MIST isochrones and mass tracks, VOSA
spectral-energy-distribution fitting, IRAF, DS9, CHARA MIRC-X/MYSTIC reduction and
interferometric model fitting}

\skill{Instrumentation \& Computing:}{Arduino, Geant4, LabVIEW, macOS, Windows, Linux}

\skill{Research Interests:}{Binary evolution, chemically peculiar stars, stellar
interferometry, stellar atmospheres, white dwarfs, astronomical data science, astrobiology,
and science communication}

\skill{Professional Affiliations:}{Sigma Pi Sigma; Society of Physics Students}

% ------------------------------------------------------- professional
\section{Professional Experience}

\entry{Vertigo Pinball, LLC -- Startup Consultant and Manager}{December 2020--July 2026}
\begin{items}
  \item Researched the barcade industry, prepared profit-and-loss projections, selected an
        arcade-machine lineup, and designed the proposed physical layout.
  \item Founded and managed the Georgia Pinball Museum initiative, curating interactive
        exhibits that served hundreds of guests daily during summers 2023 and 2024.
  \item Led Blue Ridge Gourmet Popcorn production, packaging, and specialty-flavor
        development from Summer 2024 through July 2026.
\end{items}

\entry{British Swim School -- Swim Instructor \& Substitute}{Spring 2022--Winter 2025}
\begin{items}
  \item Taught private and group lessons for students ranging from infants to older adults,
        emphasizing water safety, technique, confidence, and incremental goal-setting.
\end{items}

\entry{Huntcliff Club -- Assistant Head Swim Coach}{Summer 2016--Summer 2021}
\begin{items}
  \item Coached swimmers ages 4--18, developed competitive lineups using TeamUnify, and
        helped build a strong team culture.
\end{items}

\entry{Camp-Run-A-Mutt}{Spring 2020--Spring 2022}
\begin{items}
  \item Cared for 5--25 dogs at a cage-free day and overnight boarding facility and
        coordinated closely with staff and pet owners.
\end{items}

% --------------------------------------------------------- leadership
\section{Leadership \& Activities}

\entry{Eagle Scout, Boy Scouts of America Troop 463}{Awarded August 2020}
\begin{items}
  \item Planned and led a landscape redesign for Mount Vernon Presbyterian Church,
        including installation of hundreds of grass bundles and three Japanese maple trees.
\end{items}

\entry{Georgia State University Perimeter SPACE Club -- Secretary}{Fall 2022--Spring 2024}
\begin{items}
  \item Coordinated elementary-school outreach using physics demonstrations and
        virtual-reality experiences of ISS experiments and operations.
\end{items}

\end{document}
